\section{Отладчик GDB}
\label{sec:gdb} \index{gdb}

\subsection {Основная информация}
\label{gdb:general}

GDB (GNU Project debugger) -- отладчик, программа, позволяющая анализировать и управлять
исполнением другой программы для нахождения ошибок в коде.
GDB поддерживает отладку программ, написанных на следующих языках:

\begin{itemize}
	\item Ada
	\item Assembly
	\item C/C++
	\item D
	\item Fortran
	\item Go
	\item Objective-C
	\item OpenCL
	\item Modula-2
	\item Pascal
	\item Rust
\end{itemize}

\subsection {Основные возможности}
\label{gdb:powers}

GDB позволяет запускать программы под своим управлением, анализировать и управлять их
исполнением.

Одна из основных возможностей это остановка исполнения программы в некоторой точке.
GDB позволяет установливать специальные метки в программе -- breakpoint-ы (точки останова),
которая указывает, в каком месте следует остановить выполнение программы. Breakpoint может
быть привязан как к строке в коде, так и к опредленному адресу. Во втором случае исполнение
программы будет остановлено перед исполнением инструкции, расположенной по указанному адресу.

После остановки программы, GDB позволяет выполнять множество различных операций,
полезных для поиска ошибок в коде.

Например, GBD позволяет просматривать значение переменных и изменять их. Для низкоуровневой
отладки существует аналогичная возможность отслеживания и редактирования значений определенных
регистров.

Следующей важной возможностью GDB является просмотр информации о текущем стековом фрейме и
стеке вызова функций.

Также, весьма полезной функцией является вызов отдельных функций программы.

\subsection {Отладочная информация: DWARF}
\label{gdb:debuginfo}

Отладочная информация -- данные, которые компилятор вставляет в исполняемые файлы. Эти
данные позволяются GBD понять, какие существуют переменные, функции, их местонахождение,
как соотносяться между собой машинные инструкции и исходный код и т.~д.

DWARF -- это соврменный формат отладочной информации, который используется в компиляторах GCC
и LLVM и поддерживается GDB.

Стандарт DWARF разработан так, чтобы быть применимым к множеству существующих языков, а также
иметь возможности к расширению для будущих языков.

\newpage
